\section{Limites e continuidade em coordenadas}

Ao estudar funções que saem de um domínio e chegam em $\mathbb R^m$, é muito frequente que possamos reduzir o estudo de suas propriedades ao estudo de suas funções coordenadas.

Lembremos a definição de função coordenada.

\begin{definition}
    Sejam $A$ e $B$ conjuntos e seja $f: A \to B^m$ uma função.

    Para cada $a \in A$, $f(a)$, sendo um elemento de $B^m$, é uma $m$-upla $(b_1, \dots, b_m)$.
    As coordenadas $b_i$ são denotadas por $f_i(a)$.

    Dessa forma, $f_1(a), \dots, f_m(a)$ são os únicos elementos de $B$ tais que
    \[
        f(a) = (f_1(a), \dots, f_m(a)).
    \]

    Para cada $i\in\{1,\dots,m\}$, a função $f_i: A \to B$ assim definida é chamada de \emph{$i$-ésima função coordenada} de $f$.
\end{definition}

\begin{example}
    Seja $f: \mathbb R^2\to \mathbb R^3$ dada por $f(x, y) = (x+y, x-y, xy)$.
    Então as funções coordenadas de $f$ são dadas por:

    \begin{align*}
         & f_1: \mathbb R^2\to \mathbb R, \quad f_1(x, y) = x+y;\\
         & f_2: \mathbb R^2\to \mathbb R, \quad f_2(x, y) = x-y;\\
         & f_3: \mathbb R^2\to \mathbb R, \quad f_3(x, y) = xy.
    \end{align*}
\end{example}

Um caso importante é o das projeções, que são as funções coordenadas da função identidade de $\mathbb R^n$.

\begin{definition}
    Para cada $i\in\{1,\dots,n\}$, a $i$-ésima projeção de $\mathbb R^n$ em $\mathbb R$, denotada por $\pi_i$, é a $i$-ésima função coordenada da função identidade de $\mathbb R^n$.

    Ou seja, $\pi_i: \mathbb R^n\to \mathbb R$ é a função dada por

    \[
        \pi_i(x_1, \dots, x_n) = x_i.
    \]
\end{definition}

\begin{lemma}
    Seja $n\geq 1$ e $i$ um inteiro com $1\leq i\leq n$.

    Então a $i$-ésima projeção de $\mathbb R^n$ em $\mathbb R$ é contínua.

    Ou seja, a função $\pi_i: \mathbb R^n\to \mathbb R$ dada por
    \[
        \pi_i(x_1, \dots, x_n) = x_i
    \]

    é contínua em todos os pontos de $\mathbb R^n$.
\end{lemma}

\begin{proof}
    Seja $p=(p_1, \dots, p_n) \in \mathbb R^n$ arbitrário e $\epsilon>0$.

    Devemos ver que existe $\delta>0$ tal que, para todo $x=(x_1, \dots, x_n) \in \mathbb R^n$, se $d(x, p) < \delta$, então $d(\pi_i(x), \pi_i(p)) < \epsilon$.

    Antes de sugerir um $\delta$, vamos tomar um $\delta>0$ arbitrário e explorar as implicações de $d(x, p) < \delta$.

    Se $x\in \mathbb R^n$ é tal que $d(x, p) < \delta$, então

    \[
        \sqrt{(x_1 - p_1)^2 + \cdots + (x_n - p_n)^2} < \delta.
    \]

    Ora, então,

    \[
        d(\pi_i(x), \pi_i(p))=|x_i - p_i| = \sqrt{(x_i - p_i)^2} \leq \sqrt{(x_1 - p_1)^2 + \cdots + (x_n - p_n)^2} < \delta.
    \]

    Assim, tomando $\delta = \epsilon$, temos que $d(\pi_i(x), \pi_i(p)) < \epsilon$ sempre que $d(x, p) < \delta$.
    Isso atesta a continuidade de $\pi_i$ em $p$.

    Como $p\in \mathbb R^n$ foi arbitrário, concluímos que $\pi_i$ é contínua em todos os pontos de $\mathbb R^n$.
\end{proof}

Note que, se $f:A\to \mathbb R^m$ é uma função, sendo $f_1, \dots, f_m$ suas funções coordenadas, temos que, para cada $i\in\{1,\dots,m\}$,
\[
    f_i = \pi_i \circ f,
\]
em que $\pi_i:\mathbb R^m\to\mathbb R$ é a $i$-ésima projeção de $\mathbb R^m$ em $\mathbb R$.

De fato, se $a\in A$, então $f_i(a)$ é a $i$-ésima coordenada de $f(a)$, ou seja,
\[
    f_i(a)=\pi_i(f(a)).
\]

\begin{proposition}\index{limites!via coordenadas}
    Sejam:
    \begin{itemize}
        \item $(M, d_M)$ um espaço métrico,
        \item $A\subseteq M$ um subespaço,
        \item $f: A \to \mathbb R^m$ uma função,
        \item $p \in M$ um ponto de acumulação de $A$, e
        \item $L=(L_1, \dots, L_m)\in \mathbb R^m$.
    \end{itemize}

    Então $L$ é limite de $f$ em $p$ se, e somente se, para todo $i \in \{1, \dots, m\}$, $L_i$ é limite de $f_i$ em $p$.
\end{proposition}

\begin{proof}
Primeiro, suponha que $L$ é limite de $f$ em $p$.
Dado $i \in \{1, \dots, m\}$, temos que $f_i=\pi_i\circ f$ e que a função $\pi_i: \mathbb R^m\to \mathbb R$ é contínua. Portanto,
\[
    \lim_{x\to p} f_i(x)
    =
    \lim_{x\to p} \pi_i(f(x))
    =
    \pi_i\left(\lim_{x\to p} f(x)\right)
    =
    \pi_i(L)
    =
    L_i.
\]

    Reciprocamente, suponha que para todo $i$, $L_i$ é limite de $f_i$ em $p$.

    Devemos ver que, dado $\epsilon>0$, existe $\delta>0$ tal que, para todo $x \in A\setminus\{p\}$, se $d_M(x, p) < \delta$, então $d(f(x), L) < \epsilon$.
    Fixe $\epsilon>0$ arbitrário.

    Antes de prosseguirmos com $\epsilon$, fixemos um $\bar \epsilon>0$ arbitrário.
    Nós iremos escolher um $\bar \epsilon$ conveniente posteriormente.

    Por hipótese, para cada $i\in \{1, \dots, m\}$, existe $\delta_i>0$ tal que, para todo $x\in A\setminus\{p\}$, se $d_M(x, p) < \delta_i$, então $|f_i(x) - L_i| < \bar \epsilon$.

    Seja $\delta = \min\{\delta_1, \ldots, \delta_m\}$. Então, se $x \in A\setminus\{p\}$ é tal que $d_M(x, p) < \delta$, temos que:

    \begin{align*}
        d(f(x), L) &= \sqrt{(f_1(x) - L_1)^2 + \cdots + (f_m(x) - L_m)^2} \\
        &< \sqrt{\bar \epsilon^2+\dots+\bar \epsilon^2}= \sqrt{m}\bar \epsilon.
    \end{align*}

    Ou seja:
    \begin{equation}\label{eq:limiteViaCoordenadas}
        d(f(x), L) < \sqrt{m}\bar \epsilon.
    \end{equation}

    Assim, apliquemos a construção acima para $\bar \epsilon = \dfrac{\epsilon}{\sqrt{m}}$, obtendo $\delta$ como antes.
    Então, pela desigualdade \eqref{eq:limiteViaCoordenadas}, se $x \in A\setminus\{p\}$ é tal que $d_M(x, p) < \delta$, temos que
    \[
        d(f(x), L) < \sqrt{m}\bar \epsilon = \sqrt{m}\frac{\epsilon}{\sqrt{m}} = \epsilon.
    \]
    Como $\epsilon>0$ foi arbitrário, concluímos que $L$ é limite de $f$ em $p$.
\end{proof}

A Figura~\ref{fig:continuidadeCoord} ilustra a ideia da demonstração acima no caso $m=2$.

\subimport{continuidadeCoord}{imageLimite}


\begin{corollary}\index{continuidade!via coordenadas}
    Sejam:
    \begin{itemize}
        \item $(M, d_M)$ um espaço métrico,
        \item $f: M \to \mathbb R^m$ uma função, e
        \item $p \in M$.
    \end{itemize}

    Sejam $f_1, \dots, f_m$ as funções coordenadas de $f$.

    Então $f$ é contínua em $p$ se, e somente se, para todo $i \in \{1, \dots, m\}$, $f_i$ é contínua em $p$.

    Consequentemente, $f$ é contínua se, e somente se, para todo $i \in \{1, \dots, m\}$, $f_i$ é contínua.
\end{corollary}

\begin{proof}
    Se $p$ é um ponto isolado de $M$, trivialmente $f$ é contínua em $p$ e cada $f_i$ é contínua em $p$, pois a continuidade é automática em pontos isolados.

    Resta considerar o caso em que $p$ é um ponto de acumulação de $M$.
    Nesse caso, pela proposição anterior:

    \begin{align*}
        & f \text{ é contínua em } p \\
        \Leftrightarrow\quad & \lim_{x\to p} f(x) = f(p) \\
        \Leftrightarrow\quad & \forall i \in \{1, \dots, m\}\, \lim_{x\to p} f_i(x) = f_i(p) \\
        \Leftrightarrow\quad & \forall i \in \{1, \dots, m\}\, f_i \text{ é contínua em } p.
    \end{align*}

    A afirmação sobre continuidade de $f$ segue aplicando a equivalência acima a cada ponto $p\in M$.
\end{proof}

\begin{example}
    Seja $f:\mathbb R\to \mathbb R^2$ dada por $f(x)=(x^2, \sin x)$.

    As funções coordenadas de $f$ são $f_1(x)=x^2$ e $f_2(x)=\sin x$.

    Do nosso conhecimento sobre funções de uma variável, sabemos que $f_1$ e $f_2$ são funções contínuas.
    Assim, $f$ é contínua.
\end{example}

\begin{example}
    Seja $g:\mathbb R\to \mathbb R^2$ dada por

    \[
        g(x) = \begin{cases}
            (x^2, \sin x), & \text{se } x\neq 0;\\
            (0, 1), & \text{se } x=0.
        \end{cases}
    \]

    As funções coordenadas de $g$ são $g_1:\mathbb R\to \mathbb R$ e $g_2:\mathbb R\to \mathbb R$ dadas por

    \[
        g_1(x) = \begin{cases}
            x^2, & \text{se } x\neq 0;\\
            0, & \text{se } x=0,
        \end{cases}
    \]

    \[
        g_2(x) = \begin{cases}
            \sin x, & \text{se } x\neq 0;\\
            1, & \text{se } x=0.
        \end{cases}
    \]

    Note que $g_1(x)=x^2$ para todo $x\in \mathbb R$, e, portanto, $g_1$ é contínua.

    Por outro lado, o conjunto de pontos de continuidade de $g_2$ é $\mathbb R\setminus\{0\}$.

    Portanto, o conjunto de pontos de continuidade de $g$ é a interseção do conjunto de pontos de continuidade de $g_1$ com o conjunto de pontos de continuidade de $g_2$, ou seja, $\mathbb R\setminus\{0\}$.
\end{example}

\begin{example}
    Seja $f:\mathbb R^2\rightarrow \mathbb R^2$ dada por

    \[
        f(x, y)=(\cos(x), \sin(y)).
    \]
    
    A função $f$ é contínua.
    De fato, as funções coordenadas de $f$ são as funções $f_1,f_2:\mathbb R^2\to \mathbb R$ dadas por
    \[
        f_1(x, y)=\cos(x), \quad f_2(x, y)=\sin(y).
    \]

    Temos que $f_1$ é contínua, pois $f_1=\cos\circ \pi_1$, composição de funções contínuas.
    Analogamente, $f_2$ é contínua, pois $f_2=\sin\circ \pi_2$, composição de funções contínuas.
    Portanto, $f$ é contínua.
\end{example}

\section*{Exercícios}

\begin{exercise}
    Em cada item abaixo, justifique a continuidade da função indicada.

    \begin{enumerate}[label=(\alph*)]
        \item $f:\mathbb R^2\to\mathbb R^3$ dada por
        \[
            f(x,y)=(x,y,0).
        \]

        \item $g:\mathbb R\to\mathbb R^2$ dada por
        \[
            g(x)=(x^2,\sin x).
        \]

        \item $h:[0,+\infty[\to\mathbb R^2$ dada por
        \[
            h(x)=(\sqrt{x},x).
        \]

        \item $\varphi:]0,+\infty[^2\to\mathbb R^3$ dada por
        \[
            \varphi(x,y)=\left(\sqrt{x},\frac{1}{y},\sin x\right).
        \]
    \end{enumerate}
\end{exercise}

\begin{exercise}
    Seja $g:\mathbb R\to\mathbb R^2$ dada por
    \[
        g(x)=
        \begin{cases}
            (x,\sin x), & \text{se } x\neq 0,\\
            (0,1), & \text{se } x=0.
        \end{cases}
    \]

    Determine o conjunto dos pontos de continuidade de $g$.
\end{exercise}

\begin{exercise}
    Seja $A=\mathbb R^2\setminus\{(1,2)\}$ e seja
    $f:A\to\mathbb R^3$ dada por
    \[
        f(x,y)=(y,x,5).
    \]

    Calcule justificando com os resultados já apresentados.
    \[
        \lim_{(x,y)\to(1,2)} f(x,y).
    \]
\end{exercise}